\documentclass[10pt,a4wide]{report}

\usepackage{amsthm,amssymb,mathrsfs,setspace,pstcol,textcomp,amsmath}%amsmath, latexsym,footmisc
%\usepackage{natbib}
%\usepackage{play}
\usepackage{epsfig}
%\usepackage[grey,times]{quotchap}
\usepackage[nottoc]{tocbibind}
\renewcommand{\chaptermark}[1]{\markboth{#1}{}}
\renewcommand{\sectionmark}[1]{\markright{\thesection\ #1}}
\usepackage[margin=0.5in]{geometry}
%\usepackage{natbib}
%\bibliographystyle{unsrtnat}
%\bibliographystyle{pjgsm}

\input xy
\xyoption{all}


\theoremstyle{plain}
\newtheorem{theorem}{Theorem}[section]
\newtheorem{lemma}[theorem]{Lemma}
\newtheorem{corollary}[theorem]{Corollary}
\newtheorem{proposition}[theorem]{Proposition}

\theoremstyle{definition}
\newtheorem{definition}[theorem]{Definition}
\newtheorem{example}[theorem]{Example}
\newtheorem{notation}[theorem]{Notation}

\theoremstyle{remark}
\newtheorem{remark}[theorem]{Remark}

\renewcommand{\baselinestretch}{1.5}




\begin{document}

%\pagenumbering{arabic} \setcounter{page}{1}

% --------------- Title page -----------------------

\begin{titlepage}
\enlargethispage{3cm}

\begin{center}

\vspace*{-2cm}

\textbf{\Large DESCRETE UNIT SQUARE COVER PROBLEM}

\vfill

 A Project Report Submitted \\
 for the course \\
\vfill
%   for the Degree of  \\[10pt]

 {\large \bf MA498 Project-I}\\[5pt]
%  in \\
%  {\large \bf Mathematics and Computing}
\vfill
{\large \emph{by}}\\[10pt]
{\large\bf {Raj Kamal}}\\
{\large (Roll No. 09012321)}\\


\vfill
\includegraphics[height=2.5cm]{1.png}

\vspace*{0.5cm}

{\em\large to the}\\[10pt]
{\bf\large DEPARTMENT OF MATHEMATICS} \\[5pt]
{\bf\large \mbox{INDIAN INSTITUTE OF TECHNOLOGY GUWAHATI}}\\[5pt]
{\bf\large GUWAHATI - 781039, INDIA}\\[10pt]
{\it\large November 2012}
\end{center}

\end{titlepage}

\clearpage

% --------------- Certificate page -----------------------
\pagenumbering{roman} \setcounter{page}{2}
\begin{center}
{\Large{\bf{CERTIFICATE}}}
\end{center}
%\thispagestyle{empty}


\noindent
This is to certify that the work contained in this project report
entitled ``{\bf Descrete unit square cover problem}'' submitted
by {\bf Raj Kamal} ({\bf Roll No.: 09012321})  to Indian 
Institute of Technology Guwahati for partial requirement of {\bf MA498 Project-I} in Mathematics 
and Computing  has been carried out by him under my supervision.
\vspace{4cm}

\noindent Guwahati - 781 039 \hfill (Dr. Gautam K. Das)

\noindent   November  2012 \hfill  Project Supervisor

\clearpage

% --------------- Abstract page -----------------------
\begin{center}
{\Large{\bf{ABSTRACT}}}
\end{center}

Given a set $A$ of $n$ red points and a set $B$ of $m$ blue points on a 2-dimensional plane, 
covering points with unit squares is about finding a subset $D$ of $B$ of minimal cardinality such
that if we draw unit squares with centres around points in $D$ all the red points from A will lie
inside the unions of squares in $D$. The purpose of this project is to come up with an efficient approximation scheme to  
cover points with unit squares. This is a variant of set cover problem. The decision version of set covering is NP-complete. Lund and  Yannakakis \cite{LY94} showed that set covering cannot be 
approximated in polynomial time within a factor of $0.72 \ln n$.
\clearpage
\tableofcontents
 \index{content}
\chapter {Introduction}
A set $A$ of red points and another set $B$ of blue points are given
in 2-dimensional plane, unit squares are drawn taking blue point as center. Now descrete unit
square cover problem is finding minimal subset $D \subseteq B$, such that all red points 
lie with in the unit squares drawn taking blue points from D. 
This is an instance covering problem. In combinatorics and computer science, covering problems are computational problems 
that ask whether a certain combinatorial structure covers another, or how large the structure 
has to be to do that. Covering problems are minimization problems. Packing problems are dual of 
covering problems. Wireless sensor networks are formed by connected sensors that each have the ability
to collect, process, and store environmental information as well as communicate with
others via inter-sensor wireless communication. These characteristics allow wireless
sensor networks to be used in a wide range of applications. Covering points with unit disk finds its application in 
designing wireless sensor networks. Another problem where covering with squares or rectangle, has an 
important application is related to image processing, discussed in Tanimoto and Fowler \cite{TF80}.
None of the above problems was reported to have bounded error ratio approximation algorithm. We
shall call an algorithm a $\delta$-approximation, $\delta >0$, for a certain problem if the error of the
value of the optimal solution does not exceed $\delta$. Our main focus is to identify a $\delta$ approximation
algorithm such that $\delta$ is as small as possible. In some case one specify a family of algorithm such that for each 
$\epsilon >0$ there is an $\epsilon$-approximation algorithm in the family that solves a given problem instance 
with in relative error $\epsilon$. Such a family is called an approximation scheme. The running time of approximation
algorithm increases monotonically with $\frac{1}{\epsilon}$. If the functional dependency of the running time is on the
size of the input and $\frac{1}{\epsilon}$ is polynomial, then the scheme is said to be fully polynomail, if 
on the other hand, it is polynomial only in the input scheme, the scheme is called polynomal.
All the problems descibed and consequently their extensions are NP-Complete. As such there is no fully
polynomial scheme for these problems, unless NP=P \cite{GMJD78}.

\chapter{Related work}
                 \begin{itemize}
                 \item \textbf{Minimum geometric disk cover} Here we are given set of points and we have to find 
                              unit disk of minimum cardinality whose union covers the points. Here disk 
                              centers are not restricted to the set but can be arbitary points from the plane.
                              Again, this problem is NP-hard \cite{FMS81} and has a PTAS solution [\cite{T91}, \cite{DW85}].
                 \item \textbf{Descrete K centres }Given two sets of point in the plane $P$ and $Q$ 
                              and an integer $K$ and we have to find a set of $K$ disks centered on points in $P$ whose union covers
                              $Q$ such that the radius of the largest disk is minimized. This is contraint optimization problem.
                              We have to get a set $Q$ that consists of $K$ disks whose centers are from points in $P$ and constraint is that
                              disks can have atmost radius one. This problem is NP-Hard.
                 \item \textbf{Descrete unit disk cover Problem} Given a set $P$ of $n$ points and a set $D$ of $m$ unit disks and DUDC problem is finding 
                              minimum cardinality subset $D^{\ast}$ covering all the points in $P$\cite{GRALB12}. The research was 
                              done by Das et al., where they proposed constant factor approximation algorithm to the problem in 
                              O(n$\log$ n + m $\log$ m $+$ mn). They brought down the approximation approximation factor to 18 which 
                              was previously 22 by Claude et al., \cite{FGRSR10}.
                 \end{itemize}
\chapter {Preliminaries}
\begin{description}

\item \textbf{Restricted line separable descrete unit disk cover}(RLSDUDC):Suppose we are given a set $P$ of $n$ points and another set $D$ of
                         $m$ points such that $\forall p_{i} \in P$ $\exists$ at least one point $d_{i}$ from $D$ which covers it i,e if we draw 
                         unit disk center around $d_{i}$, then $p_{i}$ will be inside the disk and points in $P$ and center of disks in $D$ are 
                         separated by line. If such kind of condition exist then there exist an optimal set $D^{\ast} \subseteq D$ such that 
                         each point in $P$ is covered by drawing unit disks with centeres at points in $D^{\ast}$.
                         RLSDUDC is a problem of finding that optimal $D^{\ast}$ \cite{FGRSR10}.
                 
\item \textbf{Line Separable Descrete Unit Disk Cover}(LSDUDC):Suppose we have set $P$ of $n$ point and a set $D$ of $m$ points in a plane.
                         Line $\ell$ divides the plane into two halves $\ell^{+}$ and $\ell^{-}$. Set $P$ lies in $\ell^{+}$.  
                         Set $D$ which contains center of $m$ unit disks lies in $\ell^{+} \cup \ell^{-}$. $\forall p_{i} \in P$ $\exists$ at 
                         least one point $d_{i}$ from $D$ which covers it i,e if we draw unit disk center around $d_{i}$, then $p_{i}$ will be 
                         inside the disk and points in $P$. LSDUDC is problem of finding minimum sized set $D^{*} \subseteq D$ that 
                         cover all points in $P$ with the constriants describe above \cite{FGRSR10}.
             \begin{theorem}
              The RLSDUDC problem can be solved optimally where authors \cite{FGRSR10} solves the LSDUDC problem with a 2-factor approximation. The running
              time for both solutions is $O(n log n + mn)$ where $m = |D|$ and $n = |P|$ .

             \end{theorem}
             \begin{theorem}
              Descrete unit disk cover problem has an 18-approximation result in $O(m \log m+ n \log n + mn)$ time \cite{GRALB12}.

             \end{theorem}


\end{description}
\chapter{Testing Feasiblity}Checking feasiblity is an important step while solving an optimization problem. In descrete unit square cover problem
                            we have to find minimal subset set of blue points from a given set of blue points in 2 dimensional plane such that unit
                            squares drawn around points in subset covers all the red points which are also given in 2 dimensional plane. So here 
                            we are searching for minimal subset.In order to proceed for the solution first we have to be sure whether the solution 
                            exists or not i.e, checking the feasiblity of problem. For testing feasiblity, we first sort m red point and n blue points in 2 
                            dimensional plane with respect to x-coordinate to get A and B respectively. For each red point $r_{i} \in A$, where A 
                            is sorted set of n red points with respect to x-coordinate, first collect blue points $b_{i} \in B$ where B is sorted
                            set of blue points with respect to x-coordinate, whose x-coordinate lies in the $[x_{r_{i}}-0.5,x_{r_{i}}+0.5]$. Suppose 
                            we get a set $B_{i}$ of blue points then for each blue point $b_{i} \in B_{i}$
                            we calculate euclidean distance between $b_{i}$and $r_{i}$, if the distance calculated is found out to be less than one
                            then $r_{i}$ can be covered by $b_{i}$. Similarly we do for each red points $r_{i} \in A$. If we are able to get 
                            atleast a blue point for each red point in A satisfying euclidean distance less than one then the descrete unit 
                            square cover problem has a solution, otherwise not. Time Complexity ($m\log m+ n\log n + mn$).
 

\chapter{Descrete unit square cover problem}
                            Descrete unit square cover problem is solvable using the descrete unit disc cover \cite{GRALB12} in 18 approximation.
                            I have tried solving problem combining shifting strategy \cite{DW85} and using the result of line separable descrete 
                            unit disc cover \cite{FGRSR10} but the approximation was 32.Here is my approach.
                            Given a set $A$ of $n$ red points and a set $B$ of $m$ blue points in two dimensional plane.
                            \begin{itemize}
                             \item $m$ red points are sorted in X and Y direction separetely and stored in $R_{x}$ and $R_{y}$ array where $R_{x}$ and 
                                   $R_{y}$ contains points in sorted X and Y coordinate respectively.
                             \item $n$ blue points are sorted in X and Y direction separetely and stored in $B_{x}$ and $B_{y}$ array where $B_{x}$ and 
                                   $B_{y}$ contains points in sorted X and Y coordinate respectively.
                             \item a $m \ast n$ bit matrix $BM$ is created and each cell is initialized to 0. For each red point $r_{i}$ from  $A$ 
                                         we consider all the points that lies in the unit square drawn taking $r_{i}$ as center. Out of those points 
                                         if we find a blue point $b_{j}$ from set B then set $BM[j][i]=1$. Thus we prepare a bit matrix $BM$ by 
                                         which we can see what all blue points $b{j} \in B$ that covers the red point $r_{i}$ by seeing the $ith$ column
                                         of BM and all those blue points corresponding to row marked 1 in the $ith$ column are points that can cover
                                         $r_{i}$.
                              \item we also create a boolean array $Rb$ whose each entry cooresponds to a $r_{i} \in A$, each entry initialized to false.
                            \end{itemize}
                            Now I partioned the 2-dimensional plane into $1*1$ square grids. Suppose $k$ grids are formed.Now I tried finding cover
                            of red points in each square grid with the blue points that are also in same grid.
                            \begin{itemize}
                             \item In each grid, we draw a line $\ell_{1}$ that divides it into two rectangles of dimension each $0.5 \ast 1$.
                                    $\ell_{0}$ and $\ell_{2}$ define the upper and lower horizontal boundary of grid. Define upper rectangle
                                   as rectangle formed by $[\ell_{0},\ell_{1}]$ as upper and lower boundary and similarly lower reactangle as rectangle formed by $[\ell_{1},\ell_{2}]$.
                                
                             \item we take each red points in upper rectangle that can be covered by squares drawn by taking blue points 
                                   in lower rectangle as center of the square and make a set $P1$ from those red points. Again we taken each 
                                   red point in lower rectangle that can be cover by square drawn taking blue point that are in upper 
                                   rectangle as center and create a set $P2$ from those red points.
			      \item Now we take $P1$ and apply LSDUDC with respect to line $\ell_{1}$ to cover points in P1. Suppose we get a set 
                                   $S1$ as the result of LSDUDC. Now same we apply for $P2$. We apply LSDUDC 
                                   with respect to line $\ell_{1}$ to cover points in $P2$, from it we get a set $S2$. Now in the 
                                   process of applying LSDUDC we mark $Rb_{i}$ as true for each red point $r_{i}$ that are covered.
                           
                             \item we create a set $cover1 = S1 \cup S2$ that corresponds to all those blue points that cover red points that are in grid.
                             \item Similarly we do the same for each grid and get $cover1,cover2,....,coverk$ and create a set $cover=cover1 \cup cover2
                                  \cup cover3 ....\cup cover k$.
                             \item Now we scan the the boolean array $Rb$ and check entries which are marked false, for those entries, scan 
                                   the BM matrix and check all those blue points that can cover it.Suppose $Rb_{i}$ is marked false which
                                   implies $r_{i}$ is not yet covered so we check the $ith$ column of $BM$ matrix and create a set $temp$
                                   containing all blue points that covers it. If out of blue points in set $temp$, there exist a blue point that is 
                                   already in cover then make that $Rb_{i}$ true, otherwise take an arbitary blue point from the $temp$ that covers 
                                   $r_{i}$ and add to set $cover$ and make $Rb_{i}$ true.
                             \item Finally we obtain a $cover$ containing blue points that can cover all red points given.
                           \end{itemize}
                            I am applying LSDUDC twice in each grid which creates 4 approximation \cite{FGRSR10}. Now for each red points in grid,it may be cover 
                            by blue points that are in left,right grids and may be also by upper and lower grids, which overall adds 8 to the approximation
                            factor. Since we are applying shifting strategy so overall approximation becomes $8\ast 4=32$ \cite{DW85} which is much higher approximation
                            than DUDC approach so the whole amount of work is not benefecial. My future work will be going to reduce the approximation factor .
                            
\clearpage
\bibliographystyle{plain}
\begin{thebibliography}{10} 

\bibitem{DW85}D. Hochbaum and W. Maass, \textit{Approximation schemes for covering and packing problems in image processing and VLSI}, 
         J. ACM, 32, (1985) 130-136.
\bibitem{FGRSR10}F. Claude, G. K. Das, R. Dorrigiv, S. Durocher, R. Fraser, A. Lopez-Ortiz, B. G. Nickerson and A. Salinger, 
                 \textit{An improved line-separable algorithm or discrete unit disk cover}, Disc. Math, Alg. and Appl. , 2, (2010) 77-87.
\bibitem{FMS81}Fowler, M. Paterson and S. Tanimoto, \textit{Optimal packing and covering in the plane are np-complete}, Inf. Proc. 
               Let. , 12, (1981) 133-137.
\bibitem{GMJD78}Garey, M. R. , and Johnson, S. \textit{Computers and Intractability: A Guide to the Theory of NP-Completeness. Freeman, San Francisco}, (1978).
\bibitem{GRALB12} Gautam K. Das, Robert Fraser, Alejandro Lopez-Ortiz, and Bradford G. Nickerson 
                \textit{On the descrete unit disk cover problem}. 5th International Workshop, WALCOM 2011, New Delhi, India, February 18-20, 2011. Proceedings
                 (2011) 146-157 .
\bibitem{LY94} Lund, Carsten, Yannakakis, Mihalis \textit{On the hardness of approximating minimization problems} J. ACM, 
41, (1994) 960-981
\bibitem{TF80} Tanimoto, S. L. , and Fowler, R. J. \textit{Covering image subsets with patches}. In Proceedings of the 5th International 
               Conference on Pattern Recognition, (1980), 835-839.
\bibitem{T91}T. Gonzalez, \textit{Covering a set of points in multidimensional space}, Inf. Proc. Let. , 40, (1991) 181-188.

\end{thebibliography}
 
\end{document}